\medskip

Pour chacune des affirmations suivantes, préciser si elle est vraie ou fausse puis justifier
la réponse donnée.

Toute réponse non argumentée ne sera pas prise en compte.

\medskip

\begin{enumerate}
\item \textbf{Affirmation 1 }: Toute suite décroissante et minorée par 0 converge vers 0.
\item On considère une suite $\left(u_n\right)$ définie sur $\N$ telle que, pour tout entier $n$,on a \linebreak $u_n \leqslant  \dfrac{-9^n + 3^n}{7^n}$.

\textbf{Affirmation 2} : $\displaystyle\lim_{n \to + \infty} u_n =- \infty$.

\item On considère la fonction suivante écrite en langage Python:
\begin{center}
\begin{tabular}{l l}
1 {\tiny $\blacktriangledown$}& def terme(N):\\
2&\qquad  U =1\\
3 {\tiny $\blacktriangledown$}&\qquad for i in range(N) :\\
4&\qquad \qquad  U=U+i\\
5&\qquad  return U\\
\end{tabular}
\end{center}


\textbf{Affirmation 3} : terme(4) renvoie la valeur 7.
\item Lors d'un concours, le gagnant a le choix entre deux prix :

\begin{itemize}
\item Prix A : il reçoit \np{1000} euros par jour pendant 15 jours ;
\item Prix B : il reçoit 1 euro le 1\up{er} jour, 2 euros le 2\up{e} jour, 4 euros le 3\up{e} jour et pendant 15 jours la somme reçue double chaque jour.
\end{itemize}

\textbf{Affirmation 4 }: La valeur du prix A est plus élevée que la valeur du prix B.

\item On considère la suite $\left(v_n\right)$ définie pour tout entier $n \geqslant 1$ par
\[v_n = \displaystyle\int_1^n \ln x\:\text{d}x.\]

\textbf{Affirmation 5 }: La suite $\left(v_n\right)$ est croissante.
\end{enumerate}
